% Bilaga D, satt ur project/register_map.md och project/spi_register_protocol.md.

\chapter{Registerkartan och SPI-protokollet}
\label{app:protocol}

Det här är kontraktet i projektets mitt, och det är skrivet för \textbf{två klasser samtidigt}:
den här kursen bygger FPGA-sidan som en SPI-slav (\chapref{c18:ch} och \ref{c19:ch}), och
parallellklassen bygger mikrokontrollersidan som master. Båda halvorna skrivs mot det här dokumentet,
inte mot varandras kod. \textbf{Är en implementation och den här bilagan oense är implementationen
fel.}

Kontraktet säger ingenting om CAN som protokoll; det förutsätter det. \canbook{} är beskrivningen,
och dess kapitel~7 är den som förklarar varför det här lagret behövs alls: en kontroller som
signalerar med encykelspulser går inte att polla över en seriebuss.

Två saker hålls isär, och det är värt att veta vilken som är vilken. \emph{Registerkartan}
(\secref{app:protocol:map} till \ref{app:protocol:order}) säger \textbf{vad} registren betyder.
\emph{Transaktionsformatet} (\secref{app:protocol:electrical} till \ref{app:protocol:timing}) säger
\textbf{hur} de nås över ledningarna.

\begin{warning}
En äldre version av det här projektet lät drivrutinen nå registren via minnesmappad I/O på
basadressen \code{0xFF200000}. Det förutsätter en processor med en adressbuss ut mot FPGA-väven,
vilket mikrokontrollern i det här systemet inte har. Registren nås över SPI, och kolumnen
\textbf{Index} nedan, inte offseten, är det som står i kommandobyten.
\end{warning}

% ------------------------------------------------------------------------------------------
\section{Registren}
\label{app:protocol:map}

\begin{booktable}{clllL}
  \thead{Index} & \thead{Register} & \thead{Offset} & \thead{Åtkomst} & \thead{Beskrivning} \\
  \midrule
  0 & \code{STATUS} & \code{0x00} & R & Bit 0: TX klar. Bit 1: RX giltig. Bit 2: Fel. \\
  1 & \code{TX_ID} & \code{0x04} & R/W & 11-bitars sändar-ID (bitarna 10--0). \\
  2 & \code{TX_DLC} & \code{0x08} & R/W & Data Length Code (bitarna 3--0, värde 0--8). \\
  3 & \code{TX_DATA_LO} & \code{0x0C} & R/W & Sändardata byte 0--3 (byte 0 = MSB). \\
  4 & \code{TX_DATA_HI} & \code{0x10} & R/W & Sändardata byte 4--7 (byte 4 = MSB). \\
  5 & \code{TX_SEND} & \code{0x14} & W & Skriv \code{0x1} för att utlösa sändning. \\
  6 & \code{RX_ID} & \code{0x18} & R & Mottagen identifierare (bitarna 10--0). \\
  7 & \code{RX_DLC} & \code{0x1C} & R & Mottagen DLC (bitarna 3--0). \\
  8 & \code{RX_DATA_LO} & \code{0x20} & R & Mottagen data byte 0--3. \\
  9 & \code{RX_DATA_HI} & \code{0x24} & R & Mottagen data byte 4--7. \\
  10 & \code{RX_ACK} & \code{0x28} & W & Skriv \code{0x1} för att kvittera och tömma
    RX-bufferten. \\
  11 & \code{ERROR_FLAGS} & \code{0x2C} & R/W & Felregistret; skriv \code{0x0} för att
    nollställa. \\
  12 & \code{TX_ABORT} & \code{0x30} & W & Skriv \code{0x1} för att tvinga tillbaka TX klar efter
    en avbruten sändning. \\
\end{booktable}

Index 13--15 är reserverade: en läsning ger \code{0x00000000}, en skrivning ignoreras.

\code{Index} är \code{Offset / 4}. Offseten finns kvar i tabellen därför att drivrutinens konstanter
är skrivna i den formen, men det är indexet som går ut på ledningen.

% ------------------------------------------------------------------------------------------
\section{\code{STATUS}}
\label{app:protocol:status}

\begin{booktable}{cLLL}
  \thead{Bit} & \thead{Betydelse} & \thead{Sätts av} & \thead{Nollställs av} \\
  \midrule
  0 & TX klar & reset; kontrollerns puls \code{tx_done}; en stigande flank på \code{error} medan
    biten är låg; en accepterad skrivning till \code{TX_ABORT} & en accepterad skrivning till
    \code{TX_SEND} \\
  1 & RX giltig & kontrollerns puls \code{rx_valid}, som också låser \code{RX_*} & en skrivning av
    \code{0x1} till \code{RX_ACK} \\
  2 & Fel & en \textbf{stigande flank} på kontrollerns nivå \code{error} & en skrivning av
    \code{0x0} till \code{ERROR_FLAGS} \\
\end{booktable}

Poängen med hela lagret ligger i den tabellen. Kontrollern signalerar med encykelspulser, 20 ns vid
50 MHz. En drivrutin som pollar över SPI ser systemet några mikrosekunder i taget i bästa fall, så en
puls är borta tusentals cykler innan nästa pollning kommer. Registerbanken gör om pulserna till
\textbf{klibbiga, pollbara nivåer med uttrycklig nollställning}.

% ------------------------------------------------------------------------------------------
\section{Semantik, register för register}
\label{app:protocol:semantics}

\begin{itemize}
  \item \textbf{\code{TX_SEND}} utlöser en enda \code{tx_req}-puls när ett värde med \textbf{bit 0
    satt} skrivs och \code{STATUS} bit 0 är satt. Övriga bitar i det skrivna ordet ignoreras. Är
    \code{STATUS} bit 0 låg ignoreras skrivningen: en sändning pågår. En skrivning med bit 0 låg gör
    ingenting. Läsning ger \code{0x00000000}, eftersom registret inte lagrar något utan \emph{gör}
    något. Skrivningar är flankhändelser på det här lagret, så en upprepad skrivning är en upprepad
    begäran, aldrig en hållen nivå.
  \item \textbf{\code{RX_ACK}} nollställer \code{STATUS} bit 1. Kräver bit 0 satt; övriga bitar
    ignoreras. Läsning ger \code{0x00000000}.
  \item \textbf{\code{TX_ABORT}} sätter \code{STATUS} bit 0 tillbaka till \code{1}, och gör
    ingenting annat. Bara en skrivning med \textbf{bit 0 satt} gör något; läsning ger
    \code{0x00000000}. Registret når inte kontrollern \textemdash{} en kontroller som avbrutit en
    sändning är redan tillbaka i vila och kan sända igen; det är bara bankens klibbiga bit som
    fastnat. Det här är nödutgången: i en riktig konstruktion ska den aldrig behövas, eftersom bit 0
    redan har två automatiska vägar tillbaka. En drivrutin som märker att den behöver den har hittat
    en bugg värd att rapportera.
  \item \textbf{\code{ERROR_FLAGS}} är en lås-och-nollställ-på-noll-konstruktion, inte ett allmänt
    register: bara en skrivning där \textbf{hela ordet är noll} nollställer. \code{0x000000FF}
    nollställer alltså inte, trots att bit 0 är satt. Bit 0 är låsningen och speglar \code{STATUS}
    bit 2; bitarna 31--1 läses som \code{0}. Namnet är i plural av historiska skäl.
  \item \textbf{Maskning vid skrivning.} \code{TX_ID} lagrar bitarna 10--0 och \code{TX_DLC}
    bitarna 3--0; överflödiga bitar läses tillbaka som noll. En skrivning av \code{0xFFFFFFFF} läses
    alltså tillbaka som \code{0x000007FF} respektive \code{0x0000000F}.
    \code{TX_DATA_LO}/\code{HI} lagrar alla 32 bitarna.
  \item \textbf{Ingen validering.} En \code{TX_DLC} på \code{0xF} lagras och lämnas vidare till
    kontrollern som den är. Att avvisa \code{dlc > 8} är drivrutinens ansvar. Att göra det på båda
    sidor är precis så de två kopiorna av regeln glider isär.
  \item \textbf{Ingen buffring på mottagarsidan.} Kommer en ny ram innan den förra kvitterats skrivs
    den över, och den nyaste vinner. Det är ett medvetet, dokumenterat val och inte en bugg; en
    riktig kontroller har en mottagningskö.
\end{itemize}

\subsection{Detaljer en drivrutin behöver}
\label{app:protocol:driver}

Punkterna nedan hör till kontraktet lika mycket som tabellerna ovan.

\begin{itemize}
  \item \textbf{Varje register är 32 bitar brett.} Alla fem bytes i en transaktion överförs alltid,
    även för ett register som bara har fyra meningsfulla bitar.
  \item \textbf{\code{STATUS} bitarna 31--3 läses som \code{0}.} Bankens läsmux nollar allt ovanför
    bit 2, så en jämförelse mot \code{0x1} är säker.
  \item \textbf{\code{ERROR_FLAGS} skiljer inte på felorsaker.} De fyra orsakerna \textemdash{}
    förlorad arbitrering, stoppfel, misslyckad CRC och mottagen fjärrram \textemdash{} ger alla samma
    enda bit.
  \item \textbf{\code{RX_DATA_LO}/\code{HI} bortom \code{RX_DLC} läses som \code{0x00}.} Kontrollern
    nollställer sin dataackumulator vid varje ramstart, så bytes ovanför den mottagna längden bär
    aldrig rester från en tidigare ram: en ram med \code{DLC = 2} som följer på en med
    \code{DLC = 8} läses som två bytes data och sex nollor. \textbf{Maskera ändå mot \code{RX_DLC}}
    \textemdash{} nollorna är kontrollerns garanti, inte ramens innehåll, och en drivrutin som
    kopierar åtta bytes rapporterar sex bytes som aldrig sändes.
  \item \textbf{\code{TX_DATA_LO}/\code{HI} bortom \code{TX_DLC} skickas aldrig}, så de behöver inte
    nollställas.
  \item \textbf{Skrivningar till \code{STATUS}, \code{RX_ID}, \code{RX_DLC} och
    \code{RX_DATA_LO}/\code{HI} ignoreras tyst.} Transaktionen fullbordas normalt; inget fel
    rapporteras.
  \item \textbf{Det finns ingen tidsgräns att härleda.} Kartan säger inte hur lång en sändning är, så
    en drivrutin som vill ge upp får välja sin timeout själv. Ett riktvärde: en maximal ram är
    omkring 130 bitar, och vid 1 Mbit/s alltså omkring 130 µs, plus väntan på en ledig buss.
  \item \textbf{En förlorad arbitrering återställer \code{STATUS} bit 0}, precis som en lyckad
    sändning gör. En drivrutin behöver alltså inte skilja på utfallen för att kunna sända igen:
    polla bit 0, och läs \code{ERROR_FLAGS} för att veta hur det gick.
\end{itemize}

\begin{aside}
\textbf{Vad ”sändningen är slut” betyder här.} Att \code{STATUS} bit 0 kommer tillbaka säger att
sändarporten är fri igen, inte att ramen levererades. En förlorad arbitrering ger tillbaka samma bit
som en ren sändning. \code{ERROR_FLAGS} är det som skiljer dem åt, och den vanliga pollningsslingan
är därför: vänta på bit 0, läs sedan \code{ERROR_FLAGS} för att avgöra om ramen ska skickas om.
\Secref{app:project:limits} listar de övriga förenklingarna mot riktig CAN.
\end{aside}

% ------------------------------------------------------------------------------------------
\section{Databyteordning}
\label{app:protocol:order}

\code{tx_data} och \code{rx_data} är 64 bitar breda, med databyte 0 i de mest signifikanta bitarna:

\begin{widedrawing}
bit 63                                                              bit 0
+--------+--------+--------+--------+--------+--------+--------+--------+
| byte 0 | byte 1 | byte 2 | byte 3 | byte 4 | byte 5 | byte 6 | byte 7 |
+--------+--------+--------+--------+--------+--------+--------+--------+
 \___________ TX_DATA_LO ___________/ \___________ TX_DATA_HI __________/
\end{widedrawing}

Namnen \code{LO} och \code{HI} syftar på registerparets ordning, inte på bitsignifikans.
\code{TX_DATA_LO} bär de \emph{första} fyra databytesen.

% ------------------------------------------------------------------------------------------
\section{Motsvarigheter i hårdvaran}
\label{app:protocol:hw}

\begin{booktable}{LL}
  \thead{Register} & \thead{Port på \code{can_controller}} \\
  \midrule
  \code{TX_ID}, \code{TX_DLC}, \code{TX_DATA_LO}/\code{HI} & \code{tx_id}, \code{tx_dlc},
    \code{tx_data} \\
  \code{TX_SEND} & \code{tx_req} (en puls) \\
  \code{STATUS} bit 0 & \code{tx_done} (en puls), låst till en nivå \\
  \code{RX_ID}, \code{RX_DLC}, \code{RX_DATA_LO}/\code{HI} & \code{rx_id}, \code{rx_dlc},
    \code{rx_data} \\
  \code{STATUS} bit 1 & \code{rx_valid} (en puls), låst till en nivå \\
  \code{STATUS} bit 2, \code{ERROR_FLAGS} & \code{error} (en nivå), låst på stigande flank \\
\end{booktable}

Översättningen mellan de två kolumnerna är exakt vad \code{register_bank} gör, och ingenting mer.

% ------------------------------------------------------------------------------------------
\section{Elektriska parametrar och ramning}
\label{app:protocol:electrical}

\begin{booktable}{lL}
  \thead{Parameter} & \thead{Värde} \\
  \midrule
  Roller & Mikrokontrollern är SPI-master, FPGA:n slav. \\
  Läge & SPI-läge 0 (CPOL = 0, CPHA = 0): SCK vilar lågt, båda sidor samplar på stigande flank. \\
  Bitordning & MSB först, i varje byte. \\
  SCK-frekvens & $\le$ 1 MHz (AVR32DB28 går på 4 MHz och delar med 4). Se
    \secref{app:protocol:timing} för varför. \\
  \code{SS} & Aktiv låg. Ramsignal: en transaktion per låg period. \\
  Spänning & AVR32DB28:s kärna går på 5 V, DE0-CV på 3,3 V, och ingen ledning passerar en
    nivåomvandlare. De fyra SPI-ledningarna sitter på \code{PORTC}, AVR DB:s andra
    matningsdomän, med \code{VDDIO2} matad från DE0-CV:s 3,3 V-skena. \\
\end{booktable}

\subsection{Pinnindelning}
\label{app:protocol:pins}

Båda sidor spikas fast här. Mikrokontrollersidan använder \code{SPI0}:s \textbf{ALT1}-pinnmux på
AVR32DB28, som lägger periferienheten på \code{PC0}--\code{PC3}; FPGA-sidan är fri, och det här
avsnittet lägger den friheten på att göra korten utbytbara.

ALT1 är inget godtyckligt val. \code{PORTC} är AVR DB:s \textbf{MVIO}-domän: den matas från
\code{VDDIO2} i stället för från \code{VDD}, så att knyta \code{VDDIO2} till DE0-CV:s 3,3 V-skena
lägger alla fyra SPI-ledningar i en 3,3 V-domän medan kärnan fortsätter gå på 5 V. Varje nivå FPGA:n
ser är då en nivå den är specificerad för, och de 3,3 V FPGA:n driver tillbaka läses mot ett
3,3 V-tröskelvärde i stället för ett på 5 V. På \code{SPI0}:s förvalda mux skulle samma fyra signaler
ligga på \code{PA4}--\code{PA7} i \code{VDD}-domänen, och bänken skulle behöva en nivåomvandlare för
att vara säker. Det är därför muxen är en del av kontraktet i stället för varje grupps eget val.

Mikrokontrollersidan är därmed skyldig två saker utöver pinnumren, och båda hör till
drivrutinsklassen: \code{PORTMUX.SPIROUTEA} måste välja \code{ALT1}, och \code{VDDIO2} måste ha
spänning innan \code{PORTC} gör någonting alls. \code{MVIO.STATUS} rapporterar om den har det.

\begin{warning}
\textbf{Det här begränsar kortet, inte bara firmware.} Hatten som bär AVR32DB28 är byggd för det här
systemet, och pinntabellen nedan ställer tre krav på dess layout. \code{PC0}--\code{PC3} bär SPI0 och
ingenting annat \textemdash{} ingen lysdiod, ingen knapp, och framför allt ingen avstudsningskondensator
på \code{SS}, som vid 1 MHz skulle göra ramflanken oanvändbar. Ingenting annat sitter någonstans på
\code{PORTC}, eftersom hela porten är en 3,3 V-domän så snart \code{VDDIO2} är det. Och \code{VDDIO2}
dras som sitt eget nät ut till en pinne, avkopplad med 100 nF nära kapseln, i stället för att knytas
till \code{VDD} på kortet. Ingen av de tre går att rätta i firmware när kortet väl är etsat.

\textbf{En enda fuse avgör om något av ovanstående är sant.} MVIO:s tvåmatningsläge väljs av
\code{MVSYSCFG} i \code{FUSE.SYSCFG1}. I enmatningsläge knyts \code{VDDIO2} internt till \code{VDD},
\code{PORTC} blir en 5 V-port, och kopplingen i det här avsnittet driver då 5 V in i FPGA-ingångar
som inte tål 5 V \textemdash{} medan bänken ser exakt likadan ut som en riktig. Bekräfta fusen mot
databladet innan den första noden kopplas, inte efter att det första kortet betett sig konstigt.
\end{warning}

\begin{booktable}{llll}
  \thead{Signal} & \thead{AVR32DB28} & \thead{DE0-CV-port} & \thead{DE0-CV-pinne} \\
  \midrule
  \code{SCK} & \code{PC2} & \code{sclk} & \code{GPIO_0(4)} \\
  \code{MOSI} & \code{PC0} & \code{mosi} & \code{GPIO_0(5)} \\
  \code{MISO} & \code{PC1} & \code{miso} & \code{GPIO_0(6)} \\
  \code{SS} & \code{PC3} & \code{ss} & \code{GPIO_0(7)} \\
  Matning I/O & \code{VDDIO2} & \textemdash{} & 3,3 V \\
  Jord & \code{GND} & \textemdash{} & \code{GND} \\
\end{booktable}

\code{GPIO_0(0)} till \code{GPIO_0(2)} är reserverade för CAN-sidan \textemdash{} \code{tx_bus},
\code{bus_en} och den wired-AND:ade busslinjen \textemdash{} och \code{GPIO_0(3)} lämnas medvetet
oanvänd som ett glapp mellan de två grupperna, så att ett kabelknippe som sitter en pinne fel inte
kortsluter något mellan domänerna.

\textbf{Varför alls spika fast dem, när Quartus inte bryr sig?} Ingenting i konstruktionen går sönder
av att en grupp väljer andra pinnar; poängen är vad som händer vid bänken. Med en tilldelning delad
av båda klasserna passar vilken mikrokontrollernod som helst i vilket FPGA-kort som helst: ett
kabelknippe kan byggas en gång i stället för per par, och en mjukvarugrupp vars partnerkort inte är
klart kan låna ett fungerande i stället för att falla tillbaka på simulering. Det sista fallet är
inte hypotetiskt \textemdash{} de två kurserna håller sin bring-up och sina slutredovisningar samma
vecka.

En grupp \emph{får} använda andra pinnar, men bara om kortet och den nod det paras med är överens, och
bara om det är nedskrivet där den andra klassen kan se det. Ett odokumenterat pinnbyte visar sig som
en nod som svarar på registerläsningar och aldrig sänder en ram, vilket är ett dyrt sätt att upptäcka
en kopplingskonvention.

% ------------------------------------------------------------------------------------------
\section{Transaktioner}
\label{app:protocol:transactions}

Varje transaktion är exakt \textbf{5 bytes} lång: en kommandobyte, sedan fyra databytes. \code{SS}
faller före kommandobyten och stiger efter den femte byten; den måste ligga låg hela transaktionen.

\begin{textdrawing}
Bit:      7    6    5    4    3    2    1    0
        +----+----+----+----+----+----+----+----+
        | W  | 0  | 0  | 0  |   register index  |
        +----+----+----+----+----+----+----+----+
\end{textdrawing}

\begin{itemize}
  \item \textbf{Bit 7 (\code{W})}: \code{1} = skrivning, \code{0} = läsning.
  \item \textbf{Bit 6--4}: reserverade, måste vara \code{0}.
  \item \textbf{Bit 3--0}: registerindexet, \code{offset / 4} ur \secref{app:protocol:map} (0--12).
\end{itemize}

\paragraph{Skrivning (\code{W = 1}).}
Mastern skickar registervärdet i de fyra databytesen, \textbf{MSB först} (bitarna 31--24 i den första
databyten). Slaven verkställer det ihopsatta värdet i det adresserade registret först när den femte
byten är hel. Vad slaven än driver på \code{MISO} under en skrivning är meningslöst; mastern kastar
det.

\paragraph{Läsning (\code{W = 0}).}
Vid slutet av kommandobyten \textbf{låser} slaven det adresserade registrets aktuella värde
\textbf{en gång}, in i sitt svarsskiftregister. De fyra databytesen klockar sedan ut det värdet på
\code{MISO}, MSB först, medan mastern skickar dummybytes \code{0x00}. Regeln om att låsa en gång är
ingen implementationsdetalj utan en del av kontraktet: de fyra bytesen tillhör alltid \emph{ett}
sammanhängande sampel av registret, taget i ett enda ögonblick \textemdash{} \code{STATUS} i
synnerhet får aldrig vara ihopsytt av två skilda ögonblick.

\paragraph{Avbrott.}
Stiger \code{SS} innan den femte byten är hel överges transaktionen \textbf{utan sidoeffekter}:
ingenting verkställs, ingen utlösare avfyras, och slaven återgår till vila, redo för nästa fallande
flank på \code{SS}. Det är det som gör slaven självläkande efter en glitch eller en master som
resettas.

\paragraph{Ogiltiga kommandon.}
En läsning av ett reserverat index (13--15) ger \code{0x00000000}; en skrivning till ett ignoreras.
En kommandobyte med någon reserverad bit (6--4) satt är ogiltig på samma sätt: ingen läsning låses
och ingen skrivning verkställs, och slaven förbrukar ändå hela fembytesramen innan den återgår till
vila. Inget av fallen är ett fel slaven rapporterar \textemdash{} det finns ingen felkanal på det här
lagret. Att upprätthålla regeln om reserverade \emph{kommando}-bitar är \code{spi_reg_bridge}s jobb;
\code{register_bank} ser aldrig mer än ett 4-bitars index.

% ------------------------------------------------------------------------------------------
\section{Tidskrav}
\label{app:protocol:timing}

\begin{itemize}
  \item Slaven använder \textbf{aldrig} \code{SCK} som klocka. Alla tre ingångarna (\code{SCK},
    \code{MOSI}, \code{SS}) synkroniseras in i 50 MHz-domänen med två vippor var, och
    \code{SCK}-flanker \emph{detekteras}, de klockas inte på (\chapref{c4:ch} för varför,
    \chapref{c18:ch} för modulen som gör det). \code{spi_slave} bygger sin synkroniserare på plats i
    stället för att instansiera \code{meta_prev}, så att \code{spi_slave} inte beror på något i
    \file{controller/}.
  \item Tvåvippssynkronisering plus flankdetektering kostar några 50 MHz-cykler per flank; vid
    SCK $\le$ 1 MHz finns det $\ge$ 50 systemcykler per SPI-bit, alltså en tiopotens marginal.
    Protokollet garanterar därför korrekt drift bara upp till 1 MHz; snabbare kan fungera, och det
    är medvetet inte utlovat.
  \item \code{MISO} ändras på \code{SCK}:s fallande flank och samplas av mastern på den stigande, som
    läge 0 föreskriver.
  \item Mastern måste höja \code{SS} mellan transaktioner. Transaktioner rygg mot rygg under en enda
    låg \code{SS} stöds inte, och att försöka är ofarligt: \textbf{en transaktion börjar bara på
    första byten efter att \code{SS} fallit}, så en sjätte byte som anländer medan \code{SS}
    fortfarande är låg kastas i stället för att tas för en ny kommandobyte.
\end{itemize}

\subsection{Sättnings- och hålltid kring \code{SS}}
\label{app:protocol:ss}

Eftersom \code{SS} synkroniseras och sedan flankdetekteras i stället för att användas direkt behöver
den några systemklockcykler för att få effekt. Vid 50 MHz är en cykel 20 ns, och talen är:

\begin{booktable}{LlL}
  \thead{Parameter} & \thead{Minimum} & \thead{Varför} \\
  \midrule
  \code{SS} låg före första stigande \code{SCK}-flanken & \textbf{60 ns} (3 cykler) & ramstartsflanken
    detekteras på tredje systemklockflanken efter att \code{SS} fallit \\
  \code{SS} hög mellan transaktioner & \textbf{60 ns} (3 cykler) & samma detektor måste se linjen
    återvända innan den kan armas igen \\
  \code{SS} låg efter sista \code{SCK}-flanken & \textbf{60 ns} (3 cykler) & så att sista byten
    registreras innan ramen rivs \\
\end{booktable}

Räkna med 100 ns för var och en i en drivrutin så är ingenting av det här nära. En AVR som växlar
\code{SS} genom \code{PORTC.OUTCLR} och \code{PORTC.OUTSET} kring en byte skriven till
\code{SPI0.DATA} är redan långt långsammare än så, så gränserna spelar roll bara om mastern driver
\code{SS} från hårdvara, eller går mycket fortare än AVR32DB28 gör.

Det finns \textbf{inget} minsta glapp mellan bytes: bytes får komma rygg mot rygg inom en
transaktion, och slaven bryr sig inte om hur länge mastern pausar mellan dem, så länge \code{SS}
ligger kvar låg.

\subsection{\code{MISO} utanför en läsning}
\label{app:protocol:miso}

\begin{itemize}
  \item \code{MISO} drivs \textbf{lågt} medan slaven är avvald, den släpps inte till högimpedans.
    Konstruktionen förutsätter en enda slav på bussen. Att lägga till en andra skulle kräva att den
    här utgången gjordes trestatad.
  \item Byten mastern klockar ut \textbf{under kommandobyten} är meningslös: slaven laddar sitt
    svarsskiftregister vid kommandobytens \emph{slut}, så det som lämnar under den är vad som råkade
    ligga kvar. Kasta den.
  \item Under en skrivning är alla fyra \code{MISO}-databytesen meningslösa också.
\end{itemize}

% ------------------------------------------------------------------------------------------
\section{Genomarbetat exempel}
\label{app:protocol:example}

Att skicka en ram (\code{id 0x123}, \code{dlc 2}, data \code{AA BB}) är fem transaktioner, och sedan
en pollning:

\begin{widedrawing}
Write TX_ID      : 81 00 00 01 23
Write TX_DLC     : 82 00 00 00 02
Write TX_DATA_LO : 83 AA BB 00 00
Write TX_DATA_HI : 84 00 00 00 00
Write TX_SEND    : 85 00 00 00 01
Read  STATUS     : 00 xx xx xx xx   -> MISO ger 00 00 00 00 medan sändningen pågår,
                                       sedan 00 00 00 01 när tx_done satt TX klar igen.
\end{widedrawing}

Lägg märke till \code{83 AA BB 00 00}: byte 0 i den mest signifikanta positionen. Det är packningen
på båda sidor om ledningen \textemdash{} drivrutinen sätter ihop den, \code{register_bank}
presenterar \code{TX_DATA_LO} som de \emph{första} fyra bytesen, och \code{can_controller} läser
\code{tx_data} med byte 0 i bitarna 63--56. Diagrammet i \secref{app:protocol:order} är bilden av
det, och \code{register_bank_tb} kontrollerar det.

Värt att läsa av samma exempel: de fyra första \code{MISO}-bytesen av en skrivning är meningslösa,
och det är \code{MISO}-byten som returneras \emph{under} vilken kommandobyte som helst också. Kasta
den.

\begin{aside}
\textbf{Att ändra i det här kontraktet.} Registerkartan och protokollet är ett kontrakt mellan två
klasser som inte kan kompilera mot varandra. Ändras något måste båda sidor veta om det innan
ändringen mergas. I praktiken: ändringar diskuteras med båda handledarna först, och dokumenten
uppdateras \emph{före} koden, aldrig efter.
\end{aside}
